\chapter{Ringraziamenti}\label{ringraziamenti}

\hfill\break

Ho inventato un paio di malattie per dare in Spin un effetto drammatico.
La Sindrome da Deperimento Cardiovascolare è un'immaginaria malattia
trasmessa dal bestiame che non ha una controparte nel mondo reale. Anche
la Sclerosi Multipla Atipica è del tutto immaginaria, ma i suoi sintomi
imitano quelli della sclerosi multipla, sfortunatamente una malattia
molto reale. Sebbene la sclerosi multipla non sia ancora curabile, sono
state presentate diverse nuove terapie promettenti, o comunque sono
all'orizzonte. I romanzi di fantascienza non dovrebbero essere confusi
con le riviste mediche, in ogni caso. Per i lettori che vogliano saperne
di più sulla sclerosi multipla, una delle migliori risorse sul Web è:

www.nationalmssociety.org.

Il futuro che ho dedotto per la gente di Sumatra e Minangkabau è
anch'esso una mia totale invenzione, ma la cultura matrilineare minang,
e la sua coesistenza con il moderno Islam, ha attirato l'attenzione
degli antropologi - vedi lo studio di Peggy Reeves Sanday, \emph{Women
	at the Center: Life in a Modern Matriarchy}.

I lettori interessati al pensiero scientifico attuale in materia di
evoluzione e futuro del sistema solare potrebbero voler dare un'occhiata
a \emph{The Life and Death of Planet Earth} di Peter D. Ward e Donald
Brownlee o \emph{Our Cosmic Origins} di Armand Delsemme per informazioni
non deformate dalla lente della fantascienza.

E ancora una volta, fra tutti coloro che mi hanno aiutato a rendere
possibile la scrittura di questo libro (e li ringrazio tutti), la palma
della migliore va a mia moglie, Sharry.
